\section[Analytic class number formula for h-]{Analytic class number formula for $\hminus$}

Let
\[
  K=\Q(\zeta_p),
  \qquad
  K^+=\Q(\zeta_p)^+,
  \qquad
  n=\frac{p-1}{2}.
\]
Then $K$ has signature $(0,n)$ and $K^+$ has signature $(n,0)$.

\subsection*{The odd $L$-product as a residue quotient}

The factorisation of the cyclotomic Dedekind zeta function into Dirichlet
$L$-functions gives
\[
  \zeta_K(s)=\prod_{\chi \bmod p} L(s,\chi),
\]
where the product runs over all characters modulo $p$. The even characters are
exactly the characters of \(\Gal(K^+/\Q)\), so similarly
\[
  \zeta_{K^+}(s)=\prod_{\substack{\chi\bmod p\\ \chi\text{ even}}} L(s,\chi).
\]
Dividing yields
\[
  \frac{\zeta_K(s)}{\zeta_{K^+}(s)}=\prod_{\chi\text{ odd}}L(s,\chi).
\]
Since every odd character modulo $p$ is non-trivial, every factor on the right
is holomorphic at $s=1$. Taking residues at $s=1$ therefore gives:

\begin{proposition}\label{prop:odd-L-product-residue}
  \lean{BernoulliRegular.residue_ready_factorization_even_odd}
  \leanok
  \uses{def:cm-real-subfield, def:dirichlet-character, def:even-odd}
One has
\[
  \prod_{\chi\text{ odd}} L(1,\chi)
  = \frac{\operatorname*{Res}_{s=1}\zeta_K(s)}{\operatorname*{Res}_{s=1}\zeta_{K^+}(s)}.
\]
\end{proposition}

\begin{proof}\leanok
The quotient \(\zeta_K(s)/\zeta_{K^+}(s)\) has a removable singularity at
$s=1$, because both numerator and denominator have simple poles there. Its
value at $s=1$ is therefore the ratio of the residues. On the other hand this
quotient is exactly the product of the odd Dirichlet $L$-functions, all of
which are regular at $s=1$.
\end{proof}

\subsection*{The relative class number formula}

The analytic class number formula gives
\[
  \operatorname*{Res}_{s=1}\zeta_F(s)
  = \frac{2^{r_1(F)}(2\pi)^{r_2(F)}h_F R_F}{w_F\sqrt{|D_F|}}
\]
for every number field $F$, where $R_F$ is the regulator, $w_F$ the number of
roots of unity, and $D_F$ the discriminant.

For the cyclotomic fields under consideration one uses the standard identities
\[
  w_K=2p,
  \qquad
  w_{K^+}=2,
  \qquad
  R_K = 2^{n-1}R_{K^+},
  \qquad
  |D_K|=p^n|D_{K^+}|.
\]
Combining these with the definition \(h^- = h_K/h_{K^+}\) gives:

\begin{proposition}\label{prop:relative-class-number-L-product}
  \lean{BernoulliRegular.hMinus_formula_via_residues}
  \lean{BernoulliRegular.hMinus_formula_of_residue_and_hPlus_cyclotomic_and_gauss}
  \leanok
  \uses{prop:odd-L-product-residue, thm:h-factorisation, def:class-number-h,
    def:class-number-hplus, def:class-number-hminus}
One has
\[
  h^- = \frac{2p\,p^{n/2}}{(2\pi)^n}\prod_{\chi\text{ odd}}L(1,\chi).
\]
\end{proposition}

\begin{proof}\leanok
Applying the analytic class number formula to $K$ and $K^+$ gives
\[
  h_K = \operatorname*{Res}_{s=1}\zeta_K(s)\cdot
  \frac{w_K\sqrt{|D_K|}}{(2\pi)^n R_K}
\]
and
\[
  h_{K^+} = \operatorname*{Res}_{s=1}\zeta_{K^+}(s)\cdot
  \frac{w_{K^+}\sqrt{|D_{K^+}|}}{2^n R_{K^+}}.
\]
Taking the quotient,
\[
  h^- =
  \frac{\operatorname*{Res}_{s=1}\zeta_K(s)}{\operatorname*{Res}_{s=1}\zeta_{K^+}(s)}
  \cdot \frac{w_K}{w_{K^+}}
  \cdot \frac{2^n R_{K^+}}{(2\pi)^n R_K}
  \cdot \sqrt{\frac{|D_K|}{|D_{K^+}|}}.
\]
Substituting the four identities above gives
\[
  h^- = \frac{2p\,p^{n/2}}{(2\pi)^n}
  \frac{\operatorname*{Res}_{s=1}\zeta_K(s)}{\operatorname*{Res}_{s=1}\zeta_{K^+}(s)}.
\]
Now apply \cref{prop:odd-L-product-residue}.
\end{proof}

\subsection*{The product of the odd Gauss sums}

For a primitive character $\chi$ modulo $p$, write
\[
  \tau(\chi)=\sum_{a\in \Z/p\Z}\chi(a)e^{2\pi i a/p}.
\]
The standard identity
\[
  \tau(\chi)\tau(\chi^{-1}) = \chi(-1)p
\]
shows that each inverse pair of odd characters contributes \(-p\) to the
product of all odd Gauss sums.

\begin{lemma}\label{lem:odd-gauss-product}
  \lean{BernoulliRegular.rawGaussProduct}
  \lean{BernoulliRegular.cyclotomicHGaussGoal_holds}
  \leanok
  \uses{def:gauss-sum, prop:gauss-norm, def:even-odd}
One has
\[
  \prod_{\chi\text{ odd}} \tau(\chi) = i^n p^{n/2}.
\]
\end{lemma}

\begin{proof}\leanok
If $p\equiv 1\pmod 4$, then the quadratic character is even, so no odd
character is self-inverse. The odd characters split into \(n/2\) inverse
pairs, each contributing \(-p\). Therefore
\[
  \prod_{\chi\text{ odd}}\tau(\chi)=(-p)^{n/2}=i^n p^{n/2}.
\]

If $p\equiv 3\pmod 4$, then the quadratic character \(\lambda\) is odd and is
the unique self-inverse odd character. The remaining odd characters split into
\((n-1)/2\) inverse pairs, so
\[
  \prod_{\chi\text{ odd}}\tau(\chi)
  = \tau(\lambda)(-p)^{(n-1)/2}.
\]
The classical quadratic Gauss sum evaluation gives
\(
  \tau(\lambda)=i\sqrt p.
\)
Hence
\[
  \prod_{\chi\text{ odd}}\tau(\chi)
  = i\sqrt p\,(-p)^{(n-1)/2}
  = i^n p^{n/2}.
\]
\end{proof}

\begin{theorem}[Analytic formula for \texorpdfstring{$\hminus$}{h-}]
  \label{thm:hminus-formula}
  \lean{BernoulliRegular.hMinus_formula}
  \leanok
  \uses{prop:relative-class-number-L-product, cor:L1-odd, lem:odd-gauss-product,
    def:gen-bernoulli}
The relative class number satisfies
\[
  h^- = 2p\prod_{\chi\text{ odd}}
  \left(-\frac12\,\Bgen{1}{\chi^{-1}}\right).
\]
\end{theorem}

\begin{proof}\leanok
There are exactly \(n=(p-1)/2\) odd characters modulo $p$. By
\cref{prop:relative-class-number-L-product},
\[
  h^- = \frac{2p\,p^{n/2}}{(2\pi)^n}
  \prod_{\chi\text{ odd}}L(1,\chi).
\]
For each odd character \(\chi\), \cref{cor:L1-odd} gives
\[
  L(1,\chi)=\frac{\pi i\,\tau(\chi)}{p}\,\Bgen{1}{\chi^{-1}}.
\]
Multiplying over all odd characters yields
\[
  \prod_{\chi\text{ odd}}L(1,\chi)
  = \left(\frac{\pi i}{p}\right)^n
    \left(\prod_{\chi\text{ odd}}\tau(\chi)\right)
    \left(\prod_{\chi\text{ odd}}\Bgen{1}{\chi^{-1}}\right).
\]
By \cref{lem:odd-gauss-product}, this becomes
\[
  \prod_{\chi\text{ odd}}L(1,\chi)
  = \frac{\pi^n i^{2n}}{p^{n/2}}
    \prod_{\chi\text{ odd}}\Bgen{1}{\chi^{-1}}
  = \frac{\pi^n(-1)^n}{p^{n/2}}
    \prod_{\chi\text{ odd}}\Bgen{1}{\chi^{-1}}.
\]
Substituting back gives
\[
  h^- = \frac{2p\,p^{n/2}}{(2\pi)^n}
        \cdot \frac{\pi^n(-1)^n}{p^{n/2}}
        \prod_{\chi\text{ odd}}\Bgen{1}{\chi^{-1}}
      = 2p\left(-\frac12\right)^n
        \prod_{\chi\text{ odd}}\Bgen{1}{\chi^{-1}}.
\]
Since there are exactly \(n\) odd characters, this is precisely
\[
  2p\prod_{\chi\text{ odd}}
  \left(-\frac12\,\Bgen{1}{\chi^{-1}}\right).
\]
\end{proof}
